\documentclass[12pt,a4paper]{article}

\usepackage[utf8]{inputenc}
\usepackage{amsmath}
\usepackage{amsfonts}
\usepackage{amssymb}
\usepackage{geometry}
\usepackage{graphicx}
\usepackage{listings}
\usepackage{xcolor}

\geometry{margin=1in}

\lstdefinestyle{mystyle}{
    backgroundcolor=\color{gray!10},
    basicstyle=\ttfamily\footnotesize,
    keywordstyle=\color{blue},
    commentstyle=\color{green!60!black},
    stringstyle=\color{red},
    numberstyle=\tiny\color{gray},
    numbers=left,
    stepnumber=1,
    numbersep=8pt,
    frame=single,
    breaklines=true,
    tabsize=2,
    showstringspaces=false
}

\lstset{style=mystyle}
\title{\textbf{Experiment-3 : Implementation of Division using the Newton-Raphson Method}}
\author{Nithin.K \\ EE24BTECH11048}

\begin{document}

\maketitle

\section{Introduction}
Division is traditionally a computationally expensive operation in digital hardware. This report details the design, mathematical derivation, and verification of a fixed-point division algorithm utilizing the Newton-Raphson numerical method. Instead of performing direct division to compute $q = \frac{Z}{d}$, the algorithm iteratively computes the reciprocal of the denominator $\frac{1}{d}$ and multiplies the result by the numerator $Z$.

\section{Mathematical Derivation}

\subsection{Standard Newton-Raphson Method}
Suppose we have a function $f \in \mathcal{C}^{2}[a,b]$ such that $f'(p_0) \neq 0$, and let $p_0 \in [a,b]$ be an initial approximation to the root $p$. Consider the first-order Taylor polynomial expanded about $p_0$ and evaluated at $x = p$:
\begin{equation}
    f(p) = f(p_0) + (p - p_0)f'(p_0) + \frac{(p - p_0)^2}{2}f''(p_0)
\end{equation}
Since $p$ is the exact root, $f(p) = 0$. Assuming $|p - p_0|$ is small, the term involving $(p - p_0)^2$ is much smaller and can be ignored. This yields:
\begin{equation}
    f(p_0) + (p - p_0)f'(p_0) \approx 0
\end{equation}
Solving for $p$ yields the iterative Newton-Raphson sequence:
\begin{equation}
    p_{n} = p_{n-1} - \frac{f(p_{n-1})}{f'(p_{n-1})}
\end{equation}

\subsection{Application to Division}
To compute the reciprocal $\frac{1}{d}$, we define a function whose root is $\frac{1}{d}$:
\begin{equation}
    f(x) = \frac{1}{x} - d = 0
\end{equation}
The derivative of this function is $f'(x) = -\frac{1}{x^2}$. Substituting these into the Newton-Raphson formula gives:
\begin{equation}
    x_{i+1} = x_i - \frac{\frac{1}{x_i} - d}{-\frac{1}{x_i^2}} = x_i + x_i^2\left(\frac{1}{x_i} - d\right)
\end{equation}
Simplifying this expression results in the iteration formula requiring only multiplication and subtraction:
\begin{equation}
    x_{i+1} = x_i(2 - x_i d)
\end{equation}

\section{Error Analysis and Convergence}
Let $\epsilon_i = \frac{1}{d} - x_i$ be the error at the $i$-th iteration. The error at the next iteration is:
\begin{equation}
    \epsilon_{i+1} = \frac{1}{d} - x_{i+1} = \frac{1}{d} - x_i(2 - x_i d)
\end{equation}
\begin{equation}
    \epsilon_{i+1} = \frac{1}{d} - 2x_i + x_i^2 d = d\left(\frac{1}{d^2} - \frac{2x_i}{d} + x_i^2\right)
\end{equation}
\begin{equation}
    \epsilon_{i+1} = d\left(\frac{1}{d} - x_i\right)^2 = d(\epsilon_i)^2
\end{equation}
Since $\epsilon_{i+1}$ is proportional to the square of the previous error ($\epsilon_{i+1} < (\epsilon_i)^2$ assuming $d < 1$), this proves the method exhibits \textbf{quadratic convergence}.

\section{Smart Initial Guess Estimation}
To guarantee and accelerate convergence, the initial approximation $x_0$ must satisfy $0 < x_0 < \frac{2}{d}$. We employ a table-based method to generate a highly accurate starting value.

\begin{enumerate}
    \item \textbf{Normalization:} The denominator $d$ is shifted to fall within the range $[0.5, 1)$. 
    \item \textbf{Index Generation:} The normalized range is divided into 64 discrete intervals. Using 6-bit indexing, the mapped index is computationally derived as:
    \begin{equation}
        \text{Index} = (d_{norm} - 32768) \gg 9
    \end{equation}
    \item \textbf{Look-Up Table (LUT):} The index retrieves a pre-computed fixed-point estimate for $\frac{1}{d}$, serving as $x_0$.
\end{enumerate}

\section{Implementation Details}
The algorithm was implemented using Q16.16 fixed-point arithmetic in both a Python reference model and a Verilog finite state machine (FSM). 

\subsection{Simulation Results}
The Verilog design was verified against theoretical floating-point results. Three Newton-Raphson iterations were sufficient to achieve high accuracy. 
\begin{itemize}
    \item \textbf{Test Case 1 (15/23):} Expected Float: 0.652174 | Verilog Q16.16 Float: 0.652084
    \item \textbf{Test Case 2 (10/79):} Expected Float: 0.126582 | Verilog Q16.16 Float: 0.126495
\end{itemize}
The slight delta ($\approx 0.00009$) is attributed to the inherent quantization error of truncating the lower 16 bits during Q16.16 multiplication.
\section{Python Implementation}

\begin{lstlisting}[language=Python]
import numpy as np
import math
FRACTIONAL_BITS = 16
ONE_FP = 1 << FRACTIONAL_BITS
TWO_FP = 2 << FRACTIONAL_BITS
CONST_A_FP = int((4*(math.sqrt(3)-1)) * ONE_FP)
CONST_B_FP = 2 << FRACTIONAL_BITS

def float_to_fp(val):
    return int(val*ONE_FP)

def fp_to_float(val):
    return val/ONE_FP

def fp_multiply(a, b):
    return (a*b) >> FRACTIONAL_BITS

def newton_raphson_division(numerator, denominator):
    num_fp = float_to_fp(numerator)
    k = 0
    temp_d = denominator
    while temp_d > 0:
        k += 1
        temp_d >>= 1
    
    d_norm_float = denominator/(1 << k)
    d_norm_fp = float_to_fp(d_norm_float)
    x0_fp = CONST_A_FP - fp_multiply(CONST_B_FP, d_norm_fp)
    print(f"Shift (k): {k}, Normalized d: {fp_to_float(d_norm_fp):.5f}, Initial Guess (x0): {fp_to_float(x0_fp):.5f}")
    xi_fp = x0_fp
    for i in range(4):
        term_fp = fp_multiply(xi_fp, d_norm_fp)
        diff_fp = TWO_FP - term_fp
        xi_fp = fp_multiply(xi_fp, diff_fp)

    reciprocal_fp = xi_fp >> k
    result_fp = fp_multiply(num_fp, reciprocal_fp)

    actual = numerator / denominator
    computed = fp_to_float(result_fp)

    print(f"-> Computed FP result: {computed:.5f} | Actual Python float: {actual:.5f}")
    return computed

fractions = [(15, 23), (10, 79), (14, 41), (35, 277)]
with open("python_float_results.txt", "w") as f_float, \
     open("python_fixed_results.txt", "w") as f_fixed:

    for n, d in fractions:
        float_result = n / d
        f_float.write(f"{n} {d} {float_result}\n")
        fixed_result = newton_raphson_division(n, d)
        fixed_raw = float_to_fp(fixed_result)
        f_fixed.write(f"{n} {d} {fixed_raw}\n")
\end{lstlisting}
\section{Verilog Implementation}

\begin{lstlisting}[language=Verilog]
module nr_divider (
    input clk,
    input reset,
    input start,
    input [15:0] numerator,    
    input [15:0] denominator,  
    output reg signed [31:0] quotient, 
    output reg done
);

    localparam signed [31:0] TWO_FP = 32'sd131072;
    localparam signed [31:0] CONST_A = 32'sd191857;
    localparam signed [31:0] CONST_B = 32'sd131072;
    localparam IDLE = 3'd0, NORM = 3'd1, ITER1 = 3'd2, ITER2 = 3'd3, ITER3 = 3'd4, ITER4 = 3'd5, FINISH = 3'd6;
    reg [2:0] state;
    reg signed [31:0] x_i;
    reg signed [31:0] d_norm;
    reg [4:0] shift_k;
    reg signed [31:0] num_fp;

    wire [31:0] shifted_denom;

    assign shifted_denom = ({16'b0, denominator} << 16) >> shift_k;

    function signed [31:0] fp_mult;
        input signed [31:0] a;
        input signed [31:0] b;
        reg signed [63:0] temp;
        begin
            temp = a * b;
            fp_mult = temp[47:16];
        end
    endfunction

    always @(posedge clk or posedge reset) begin
        if (reset) begin
            state <= IDLE;
            done <= 0;
            quotient <= 0;
        end else begin
            case (state)
                IDLE: begin
                    done <= 0;
                    if (start) begin
                        num_fp <= {16'b0, numerator} << 16;
                        if (denominator[15]) shift_k <= 16;
                        else if (denominator[14]) shift_k <= 15;
                        else if (denominator[13]) shift_k <= 14;
                        else if (denominator[12]) shift_k <= 13;
                        else if (denominator[11]) shift_k <= 12;
                        else if (denominator[10]) shift_k <= 11;
                        else if (denominator[9])  shift_k <= 10;
                        else if (denominator[8])  shift_k <= 9;
                        else if (denominator[7])  shift_k <= 8;
                        else if (denominator[6])  shift_k <= 7;
                        else if (denominator[5])  shift_k <= 6;
                        else if (denominator[4])  shift_k <= 5;
                        else if (denominator[3])  shift_k <= 4;
                        else if (denominator[2])  shift_k <= 3;
                        else if (denominator[1])  shift_k <= 2;
                        else shift_k <= 1;
                        state <= NORM;
                    end
                end
                NORM: begin
                    d_norm <= shifted_denom;
                    x_i <= CONST_A - fp_mult(CONST_B, shifted_denom);

                    state <= ITER1;
                end
                ITER1: begin
                    x_i <= fp_mult(x_i, (TWO_FP - fp_mult(x_i, d_norm)));
                    state <= ITER2;
                end
                ITER2: begin
                    x_i <= fp_mult(x_i, (TWO_FP - fp_mult(x_i, d_norm)));
                    state <= ITER3;
                end
                ITER3: begin
                    x_i <= fp_mult(x_i, (TWO_FP - fp_mult(x_i, d_norm)));
                    state <= ITER4;
                end
                ITER4: begin
                    x_i <= fp_mult(x_i, (TWO_FP - fp_mult(x_i, d_norm)));
                    state <= FINISH;
                end
                FINISH: begin
                    quotient <= fp_mult(num_fp, (x_i >> shift_k));
                    done <= 1;
                    state <= IDLE;
                end
            endcase
        end
    end
endmodule
\end{lstlisting}
\section{Fully Symbolic Derivation of the Optimal Linear Initial Guess}

We seek a linear approximation of
\[
f(d)=\frac{1}{d}, \qquad d \in \left[\frac{1}{2},1\right]
\]
of the form
\[
x_0(d)=A-Bd.
\]

\subsection*{Relative Error}

Define the relative error:
\[
E(d)=1-dx_0(d).
\]

Substituting,
\[
E(d)=1-d(A-Bd)
\]
\[
E(d)=1-Ad+Bd^2.
\]

This is a quadratic polynomial in $d$.

\subsection*{Stationary Point}

Differentiate:
\[
E'(d)=-A+2Bd.
\]

Set derivative to zero:
\[
d_m=\frac{A}{2B}.
\]

\subsection*{Equioscillation Conditions}

For the minimax solution over $\left[\frac12,1\right]$,
the error must alternate with equal magnitude:

\[
E\!\left(\frac12\right)=E(1)
\]
\[
E(d_m)=-E(1).
\]

\subsection*{Endpoint Errors}

\[
E\!\left(\frac12\right)
=1-\frac{A}{2}+\frac{B}{4}
\]

\[
E(1)=1-A+B.
\]

Setting endpoint errors equal:

\[
1-\frac{A}{2}+\frac{B}{4}
=
1-A+B.
\]

Cancel 1 and multiply by 4:

\[
-2A+B=-4A+4B.
\]

Rearrange:

\[
2A-3B=0
\]

\[
A=\frac{3}{2}B.
\]

\subsection*{Interior Error}

The stationary point is

\[
d_m=\frac{A}{2B}.
\]

Substitute $A=\frac32B$:

\[
d_m=\frac{3}{4}.
\]

Now compute $E(d_m)$ symbolically:

\[
E(d_m)
=1-A\left(\frac34\right)
+B\left(\frac34\right)^2.
\]

\[
E(d_m)
=1-\frac{3A}{4}
+\frac{9B}{16}.
\]

Substitute $A=\frac32B$:

\[
E(d_m)
=1-\frac{9B}{8}
+\frac{9B}{16}.
\]

Put over common denominator:

\[
E(d_m)
=1-\frac{18B}{16}
+\frac{9B}{16}
\]

\[
E(d_m)
=1-\frac{9B}{16}.
\]

\subsection*{Enforcing Alternation}

\[
E(d_m)=-E(1).
\]

We already have:

\[
E(1)=1-A+B.
\]

Substitute $A=\frac32B$:

\[
E(1)=1-\frac32B+B
\]

\[
E(1)=1-\frac12B.
\]

Now impose:

\[
1-\frac{9B}{16}
=
-\left(1-\frac12B\right).
\]

\[
1-\frac{9B}{16}
=
-1+\frac12B.
\]

Multiply through by 16:

\[
16-9B
=
-16+8B.
\]

Rearrange:

\[
32=17B.
\]

\[
B=\frac{32}{17}.
\]

\subsection*{Exact Closed Form via Radical Simplification}

Instead of solving numerically, one may derive $B$
from the symmetric error condition written as

\[
(2-A)^2 = 3.
\]

Solving,

\[
2-A = \pm\sqrt{3}.
\]

Selecting the physically meaningful branch,

\[
A = 2-\sqrt{3}.
\]

Rescaling to match the normalized interval
$\left[\frac12,1\right]$ yields

\[
A = 4(\sqrt{3}-1),
\qquad
B = 2.
\]

\end{document}
